% Generated from locale/tr/content/lambda-calculus/syntax/beta.tex; do not hand-edit.


\olfileid[tr]{lam}{int}{bet}
\olsection{$\beta$-indirgeme}

$(\lambd[m][(\lambd[y][y]) m])$ terimini gördüğümüzde bunun
$\lambd[m][m]$ ile bağlantılı olduğunu, ikinci terimin birincinin
``sadeleştirilmesiyle'' elde edildiğini düşünmek doğaldır.
\emph{$\beta$-indirgeme} kavramı bu sezgiyi biçimsel olarak karşılar.

\begin{defn}[$\beta$-daraltma, $\bredone$] \ollabel{defn:betacontr}
  \emph{$\beta$-daraltma} ($\bredone$), terimler üzerinde aşağıdaki koşulu
  sağlayan en küçük bağdaşır bağıntıdır:
  \[
    (\lambd[x][N])Q \bredone \Subst{N}{Q}{x}
  \]
  $P\bredone Q$ ise $P$'nin $Q$'ya \emph{$\beta$-daraltıldığını} söyleriz.
  $(\lambd[x][N])Q$ biçimindeki terime \emph{indirgenebilir ifade} ya da
  \emph{redex} denir.
\end{defn}

\begin{prob} \ollabel{prob:def}
  Bağlı değişkeni değiştirme için
  \olref[lam][syn][alp]{defn:aconvone} içinde yaptığımız gibi,
  $\beta$-daraltmanın eşdeğer tümevarımlı tanımını açıkça yazın.
\end{prob}

\begin{defn}[$\beta$-indirgeme, $\bred$] \ollabel{defn:betared}
  \emph{$\beta$-indirgeme} ($\bred$), $\bredone$ bağıntısını içeren,
  terimler üzerindeki en küçük yansımalı ve geçişli bağıntıdır. $P\bred Q$
  ise $P$'nin $Q$'ya $\beta$-indirgenmiş olduğunu söyleriz.
\end{defn}

Bağlam açık olduğunda $\bredone$ yerine $\redone$, $\bred$ yerine $\red$
yazacağız.

Gayriresmî olarak $M\bred N$ ancak ve ancak $M$, sıfır ya da daha çok
$\beta$-daraltma adımıyla $N$'ye dönüştürülebiliyorsa geçerlidir.

\begin{defn}[$\beta$-normal]
Daha fazla $\beta$-daraltılamayan terime \emph{$\beta$-normal} denir.
\end{defn}

$M\bred N$ ve $N$ $\beta$-normalse $N$'ye $M$'nin bir \emph{normal biçimi}
denir. Bir terimin normal biçiminin tek olup olmadığı sorulabilir; ileride
göreceğimiz gibi yanıt evettir.

Bazı örnekleri ele alalım.
\begin{enumerate}
\item Şunlar geçerlidir:
\begin{align*}
(\lambd[x][xxy]) \lambd[z][z] & \redone (\lambd[z][z])(\lambd[z][z]) y \\
& \redone (\lambd[z][z]) y \\
& \redone y.
\end{align*}
\item Bir terimi ``sadeleştirmek'' onu gerçekte daha karmaşık kılabilir:
\begin{align*}
(\lambd[x][xxy])(\lambd[x][xxy]) & \redone (\lambd[x][xxy])(\lambd[x][xxy])y \\
& \redone (\lambd[x][xxy])(\lambd[x][xxy])yy \\
& \redone \dots
\end{align*}
\item Bir terimi değişmeden de bırakabilir:
\[
(\lambd[x][xx])(\lambd[x][xx]) \redone (\lambd[x][xx])(\lambd[x][xx]).
\]
\item Bazı terimler birden çok yoldan indirgenebilir. Örneğin en dıştaki
uygulama daraltılırsa
\[
(\lambd[x][(\lambd[y][yx]) z]) v \redone (\lambd[y][yv]) z;
\]
en içteki daraltılırsa
\[
(\lambd[x][(\lambd[y][yx]) z]) v \redone (\lambd[x][zx]) v
\]
elde edilir. Bununla birlikte bu durumda her iki terimin de aynı $zv$
terimine indirgenmeye devam ettiğine dikkat edin.
\end{enumerate}

Son örnekteki ortak sonuç rastlantı değildir; lambda hesabının
Church--Rosser özelliği olarak bilinen derin ve önemli bir niteliğini gösterir.

\begin{digress}
  Genel olarak bir terimi $\beta$-indirgemek için birden çok yol vardır; bu
  nedenle çeşitli indirgeme stratejileri geliştirilmiştir. En yaygını
  \emph{doğal stratejidir}. Doğal strateji her zaman konumu terimdeki başlangıç
  noktasıyla belirlenen \emph{en soldaki} redexi daraltır. Bu stratejinin
  yararlı özelliği şudur: bir terim herhangi bir stratejiyle normal biçime
  indirgenebiliyorsa, ancak ve ancak doğal stratejiyle normal biçime
  indirgenebilir. Aksi belirtilmedikçe ileride doğal stratejiyi kullanacağız.
\end{digress}

\begin{defn}[$\beta$-denkliği, $\equal$]
  \emph{$\beta$-denkliği} ($\equal$) aşağıdaki gibi tümevarımlı tanımlanan
  bağıntıdır:
  \begin{enumerate}
  \item $M \equal M$.
  \item $M \equal N$ ise $N \equal M$.
  \item $M \equal N$, $N \equal O$ ise $M \equal O$.
  \item $M \equal N$ ise $PM \equal PN$.
  \item $M \equal N$ ise $MQ \equal NQ$.
  \item $M \equal N$ ise $\lambd[x][M] \equal \lambd[x][N]$.
  \item $(\lambd[x][N])Q \equal \Subst{N}{Q}{x}$.
  \end{enumerate}
\end{defn}

İlk üç kural bağıntıyı bir denklik bağıntısı, sonraki üçü bağdaşır kılar; son
kural ise bağıntının $\beta$-daraltmayı içermesini sağlar.

Gayriresmî olarak iki terim, biri diğerine sıfır ya da daha çok
$\beta$-daraltma veya ``ters $\beta$-daraltma'' adımıyla dönüştürülebiliyorsa
ve ancak o zaman $\beta$-denktir. Ters $\beta$-daraltma, $N$, $M$'ye
$\beta$-daralıyorsa ve ancak o zaman $M$'nin $N$'ye ters
$\beta$-daralacağı biçimde tanımlanır.

Yukarıdaki kurallara başka kurallar da ekleyeceğiz ve genişletilmiş denklik
bağıntısını, $X$ eklenen kural olmak üzere, $\equal[X]$ ile göstereceğiz.

